\documentclass[10pt,landscape,a4paper]{article}
\usepackage[utf8]{inputenc}
\usepackage[T1]{fontenc}
\usepackage{geometry}
\usepackage{multicol}
\usepackage{amsmath,amssymb}
\usepackage{enumitem}
\usepackage{graphicx}
\usepackage{xcolor}
\usepackage{listings}
\usepackage{booktabs}
\usepackage{titlesec}
\usepackage{fancyhdr}
\usepackage{hyperref}
\usepackage{CJKutf8}

% Page layout
\geometry{top=0.8cm, bottom=0.8cm, left=0.8cm, right=0.8cm}
\pagestyle{empty}
\setlength{\parindent}{0pt}
\setlength{\parskip}{2pt}
\setlength{\columnsep}{8pt}

% Section formatting
\titleformat{\section}{\normalsize\bfseries\color{blue!70!black}}{}{0em}{}[\titlerule]
\titleformat{\subsection}{\small\bfseries\color{blue!50!black}}{}{0em}{}
\titlespacing{\section}{0pt}{4pt}{2pt}
\titlespacing{\subsection}{0pt}{3pt}{1pt}

% Code listing style
\lstset{
    language=Python,
    basicstyle=\ttfamily\scriptsize,
    keywordstyle=\color{blue!70!black}\bfseries,
    stringstyle=\color{red!60!black},
    commentstyle=\color{green!50!black}\itshape,
    numberstyle=\tiny\color{gray},
    backgroundcolor=\color{gray!8},
    frame=single,
    framerule=0.3pt,
    rulecolor=\color{gray!50},
    breaklines=true,
    breakatwhitespace=false,
    tabsize=2,
    showstringspaces=false,
    aboveskip=3pt,
    belowskip=3pt,
    xleftmargin=2pt,
    xrightmargin=2pt,
}

% Custom box for key concepts
\newcommand{\keybox}[1]{%
    \begin{center}
    \fcolorbox{blue!50!black}{blue!5}{%
        \parbox{0.95\linewidth}{\small #1}%
    }
    \end{center}
}

% Chinese hint style (gray, smaller)
\newcommand{\zhint}[1]{{\scriptsize\color{gray!70!black}#1}}

\begin{document}
\begin{CJK}{UTF8}{gbsn}

\begin{center}
    {\large\bfseries IERG4320/IEMS5726 --- Chapter 3: Data Representation I --- Cheat Sheet}\\[2pt]
    {\small Prof.\ YAN Wenjing, Dr.\ Danny Yip \quad|\quad Department of Information Engineering, CUHK}
\end{center}

\vspace{2pt}
\begin{multicols}{3}

% ============================================================
\section{Data Preprocessing \zhint{(数据预处理)}}
% ============================================================

\zhint{在使用数据之前先把数据整理干净，让模型能更好地学习。}

\subsection{Why? \zhint{(为什么要预处理？)}}
\begin{itemize}[leftmargin=*, nosep, itemsep=0pt]
    \item Accuracy \zhint{(准确性：数据对不对？)}
    \item Completeness \zhint{(完整性：有没有缺失？)}
    \item Consistency \zhint{(一致性：各处数据是否矛盾？)}
    \item Timeliness \zhint{(时效性：数据是否过时？)}
    \item Believability \zhint{(可信度：数据可靠吗？)}
    \item Interpretability \zhint{(可理解性：能不能看懂？)}
\end{itemize}

\subsection{4 Goals \zhint{(四大目标)}}
\begin{enumerate}[leftmargin=*, nosep, itemsep=0pt]
    \item \textbf{Data Cleaning} \zhint{(清洗)}: fix/remove wrong, duplicate, incomplete data
    \item \textbf{Data Integration} \zhint{(整合)}: merge data from multiple sources
    \item \textbf{Data Reduction} \zhint{(降维)}: reduce data size while keeping info
    \item \textbf{Data Transformation} \zhint{(转换)}: change format/structure/values
\end{enumerate}

\subsection{Data Wrangling Steps \zhint{(数据整理六步)}}
\begin{enumerate}[leftmargin=*, nosep, itemsep=0pt]
    \item Discovering \zhint{(了解数据长什么样)}
    \item Structuring \zhint{(组织原始数据)}
    \item Cleaning \zhint{(统一格式、去掉异常值)}
    \item Enriching \zhint{(能不能补充更多数据？)}
    \item Validating \zhint{(检查质量和一致性)}
    \item Publishing \zhint{(准备好给下游使用)}
\end{enumerate}

\subsection{Preprocessing Pipeline \zhint{(预处理流程)}}
\keybox{
\zhint{加载数据 $\to$ 数据整理 $\to$ 处理缺失值 $\to$ 处理分类变量 $\to$ 去重 $\to$ 划分训练/验证/测试集 $\to$ 归一化 $\to$ (可选)特征选择}
}

% ============================================================
\section{Train/Val/Test Split \zhint{(数据集划分)}}
% ============================================================

\begin{itemize}[leftmargin=*, nosep, itemsep=1pt]
    \item \textbf{Training set} \zhint{(训练集)}: train the model \zhint{(用来训练模型)}
    \item \textbf{Validation set} \zhint{(验证集)}: tune hyperparams \zhint{(调参、选模型)}
    \item \textbf{Testing set} \zhint{(测试集)}: final evaluation \zhint{(最终评估，只用一次)}
\end{itemize}

\begin{lstlisting}
from sklearn.model_selection import train_test_split
# First split: 80% train+val, 20% test
X_1, X_test, y_1, y_test = train_test_split(
    X, y, test_size=0.2, random_state=0)
# Second split: 80% train, 20% val
X_train, X_val, y_train, y_val = train_test_split(
    X_1, y_1, test_size=0.2, random_state=0)
\end{lstlisting}

\subsection{Cross Validation \zhint{(交叉验证)}}
\zhint{把数据分成k份，轮流用其中1份做验证，其余做训练，重复k次取平均。}\\
\zhint{目的：更可靠地评估模型效果，避免"运气好"。}

\textbf{Model selection example} \zhint{(选模型技巧)}:\\
\zhint{选CV和Test表现都不错且差距小的模型（如M3/M4），不选差距太大的（如M1/M2）。}

% ============================================================
\section{Data Imputation \zhint{(缺失值填补)}}
% ============================================================

\zhint{数据缺失很常见（漏填、拒绝回答、设备故障等），不能直接扔掉，要想办法填上。}

\subsection{Types of Missingness \zhint{(缺失类型)}}
\begin{itemize}[leftmargin=*, nosep, itemsep=1pt]
    \item \textbf{MCAR} \zhint{(完全随机缺失)}: missing by pure chance\\
    \zhint{像随机在数据上挖洞，跟数据本身无关}
    \item \textbf{MAR} \zhint{(随机缺失)}: depends on other observed variables\\
    \zhint{缺失跟其他已知变量有关，如性别影响是否回答某问题}
    \item \textbf{MNAR} \zhint{(非随机缺失)}: depends on the missing value itself\\
    \zhint{缺失跟缺的值本身有关，如作弊的人不回答"你作弊过吗"}
\end{itemize}

\subsection{Method 1: Deletion \zhint{(删除法)}}
\begin{lstlisting}
# Drop rows by index
df.drop([0, 1])
# Drop columns by name
df.drop(columns=['A','B'], axis=1)
# Drop any row with NaN
df.dropna()
\end{lstlisting}
\zhint{简单粗暴，但会丢失信息。}

\subsection{Method 2: Mean/Median/Mode \zhint{(均值/中位数/众数填补)}}
\begin{lstlisting}
import numpy as np
from scipy import stats
# Fill with mean
df.replace({'D': np.nan}, np.mean(df['D']))
# Fill with median
df.replace({'B': np.nan},
    np.nanmedian(df['B']))
# Fill with mode
df.replace({'D': np.nan},
    stats.mode(df['D'])[0][0])
\end{lstlisting}

\subsection{Method 3: Regression \zhint{(回归填补)}}
\zhint{用其他列建回归模型来预测缺失值。}
\begin{lstlisting}
from sklearn import linear_model
df_train = df.dropna()
X = df_train.drop(columns=['D'])
regr = linear_model.LinearRegression()
regr.fit(X, df_train['D'])
regr.predict(df.drop(columns=['D']))
\end{lstlisting}
\zhint{注意：可能引入变量间额外依赖关系。}

\subsection{Method 4: Hot Deck \zhint{(热卡填补)}}
\zhint{用同一列中最近的有效值来填充（向前填充）。}
\begin{lstlisting}
# Forward fill: use last valid value
df.ffill(inplace=True)
\end{lstlisting}

\subsection{Method 5: Multiple Imputation \zhint{(多重填补)}}
\zhint{最精确的方法：建多个模型，综合分析不确定性。}
\begin{lstlisting}
from sklearn.experimental import (
    enable_iterative_imputer)
from sklearn.impute import IterativeImputer
imp = IterativeImputer(
    max_iter=10, random_state=0)
imp.fit([[1,2],[3,6],[4,8],
         [np.nan,3],[7,np.nan]])
X_test = [[np.nan,2],[6,np.nan],[np.nan,6]]
print(np.round(imp.transform(X_test)))
# [[ 1.  2.] [ 6. 12.] [ 3.  6.]]
\end{lstlisting}
\zhint{模型学到第二列是第一列的2倍，自动补全。}

\subsection{sklearn Imputers \zhint{(sklearn填补工具)}}
\begin{itemize}[leftmargin=*, nosep, itemsep=0pt]
    \item \texttt{SimpleImputer} \zhint{(均值/中位数/众数/固定值)}
    \item \texttt{IterativeImputer} \zhint{(迭代回归，多变量)}
    \item \texttt{KNNImputer} \zhint{(用K近邻来填补)}
    \item \texttt{MissingIndicator} \zhint{(标记哪里缺失)}
\end{itemize}

\subsection{Guidelines (Harrell 2001) \zhint{(实用指南)}}
\begin{center}
\scriptsize
\begin{tabular}{lp{4.5cm}}
\toprule
\textbf{Missing \%} & \textbf{Strategy} \zhint{(怎么处理)} \\
\midrule
$\leq$ 5\% & Median/Mode \zhint{(简单填就行)} \\
5--15\% & Multiple imputation if correlated \zhint{(看情况)} \\
15--30\% & Multiple imputation \zhint{(必须用多重填补)} \\
$>$ 30\% & Consider dropping the variable \zhint{(缺太多，考虑删)} \\
\bottomrule
\end{tabular}
\end{center}

% ============================================================
\section{Categorical Variables \zhint{(分类变量处理)}}
% ============================================================

\zhint{机器学习需要数字输入，所以要把文字类别转成数字。}

\subsection{Label Encoding \zhint{(标签编码)}}
\zhint{把每个类别映射成一个数字。如 male$\to$1, female$\to$0}
\begin{lstlisting}
from sklearn import preprocessing
le = preprocessing.LabelEncoder()
df['A'] = le.fit_transform(df['A'])
df['A'].unique()
\end{lstlisting}

\subsection{Binning \zhint{(分箱/分桶)}}
\zhint{把连续值分成几个区间。如 age: 0--17$\to$0, 18--64$\to$1, 65+$\to$2}
\begin{lstlisting}
bins = [0, 17, 64, np.inf]
names = [0, 1, 2]
df['ageRange'] = pd.cut(
    df['age'], bins, labels=names)
\end{lstlisting}

\subsection{Ordinal Encoding \zhint{(有序编码)}}
\zhint{类别有大小关系时使用。如 red$\to$1, green$\to$2, blue$\to$3}
\begin{lstlisting}
df.name = pd.Categorical(df['name'])
df['code'] = df.name.cat.codes
\end{lstlisting}

\subsection{One-Hot Encoding \zhint{(独热编码)}}
\zhint{每个类别变成一列0/1。red$\to$\{1,0,0\}, green$\to$\{0,1,0\}, blue$\to$\{0,0,1\}}\\
\textbf{Dummy encoding} \zhint{(哑变量)}: drop one column to avoid redundancy.\\
\zhint{red$\to$\{1,0\}, green$\to$\{0,1\}, blue$\to$\{0,0\}（少一列，减少冗余）}

\subsection{Combine Levels \zhint{(合并类别)}}
\zhint{频率很低的类别可以合并到一起，减少稀疏性。}

% ============================================================
\section{Duplication Removal \zhint{(去重)}}
% ============================================================

\subsection{Drop Duplicate Rows \zhint{(删重复行)}}
\begin{lstlisting}
data = df.drop_duplicates()
\end{lstlisting}

\subsection{Filter Constant Columns \zhint{(删常数列)}}
\zhint{如果一列所有值都一样，没有信息量，删掉。}
\begin{lstlisting}
data = df.loc[:,
    (df != df.iloc[0]).any()]
\end{lstlisting}

\subsection{Filter Correlated Columns \zhint{(删高相关列)}}
\zhint{两列高度相关（如温度和温度$\times$2），留一个就够了。}
\begin{lstlisting}
def correlation(dataset, threshold):
    col_corr = set()
    corr_matrix = dataset.corr()
    for i in range(len(corr_matrix.columns)):
        for j in range(i):
            if abs(corr_matrix.iloc[i,j]) \
                    > threshold:
                colname = \
                    corr_matrix.columns[i]
                col_corr.add(colname)
    return col_corr
\end{lstlisting}

% ============================================================
\section{Data Normalization \zhint{(数据归一化)}}
% ============================================================

\zhint{不同特征量纲差距大（如年龄1--100 vs 收入10000--1000000），需要统一尺度。}

\subsection{1. Min-Max \zhint{(最小最大归一化)}}
\zhint{缩放到[0, 1]区间：$x' = \frac{x - \min}{\max - \min}$}
\begin{lstlisting}
from sklearn import preprocessing
scaler = preprocessing.MinMaxScaler()
X_scaled = scaler.fit_transform(X_train)
\end{lstlisting}

\subsection{2. Mean Normalization \zhint{(均值归一化)}}
\zhint{$x' = \frac{x - \text{mean}}{\max - \min}$，结果在[-1, 1]附近}
\begin{lstlisting}
df_means = df.mean(axis=0)
min_max_range = df.max(0) - df.min(0)
df_scaled = (df - df_means) / min_max_range
\end{lstlisting}

\subsection{3. Z-score / Standardization \zhint{(标准化)}}
\zhint{$x' = \frac{x - \mu}{\sigma}$，均值变0，标准差变1。最常用！}
\begin{lstlisting}
from sklearn.preprocessing import StandardScaler
ss = StandardScaler()
df_scaled = ss.fit_transform(df)
\end{lstlisting}

\subsection{4. Unit Length \zhint{(单位长度缩放)}}
\zhint{让每个特征向量的长度（L2范数）变成1}
\begin{lstlisting}
df_scaled = df / np.sqrt(
    np.square(df).sum(axis=0))
\end{lstlisting}

% ============================================================
\section{Feature Selection \zhint{(特征选择)}}
% ============================================================

\zhint{数据太多列怎么办？挑出最有用的列，扔掉没用的，让模型更快更准。}

\subsection{Why? \zhint{(为什么要特征选择？)}}
\begin{itemize}[leftmargin=*, nosep, itemsep=0pt]
    \item Remove irrelevant features \zhint{(去掉没用的列)}
    \item Improve accuracy \zhint{(提高准确率)}
    \item Reduce cost \zhint{(省计算资源)}
    \item Simplify model \zhint{(模型更简单易懂)}
\end{itemize}

\subsection{Two Approaches \zhint{(两大类方法)}}
\begin{enumerate}[leftmargin=*, nosep, itemsep=1pt]
    \item \textbf{Select columns/rows} \zhint{(直接选列/行)}\\
    \zhint{不改变原始特征，只挑子集}
    \item \textbf{Dimensionality reduction} \zhint{(降维)}\\
    \zhint{把多个特征压缩成少数几个新特征（PCA/LDA）}
\end{enumerate}

\subsection{Search Directions \zhint{(搜索方向)}}
\begin{itemize}[leftmargin=*, nosep, itemsep=0pt]
    \item \textbf{SFG} \zhint{(前向)}: start empty, add features one by one
    \item \textbf{SBG} \zhint{(后向)}: start full, remove features one by one
    \item \textbf{BG} \zhint{(双向)}: add and remove simultaneously
    \item \textbf{RG} \zhint{(随机)}: random subsets
\end{itemize}

\subsection{Selection Criteria \zhint{(评价标准)}}
\begin{itemize}[leftmargin=*, nosep, itemsep=0pt]
    \item Information gain \zhint{(信息增益)}
    \item Distance measure \zhint{(距离度量)}
    \item Correlation (Pearson) \zhint{(相关系数)}
    \item Consistency \zhint{(一致性)}
    \item Accuracy \zhint{(准确率)}
\end{itemize}

\subsection{Wrapper Method \zhint{(包裹法)}}
\zhint{用分类器本身来评估哪些特征组合最好。}
\begin{lstlisting}
from mlxtend.feature_selection import (
    SequentialFeatureSelector as SFS)
sfs = SFS(RandomForestClassifier(n_jobs=8),
    k_features=5,   # pick 5 features
    forward=True,    # forward search
    floating=False,
    scoring='roc_auc', cv=3)
sfs = sfs.fit(X_train.values, y_train)
selected = X_train.columns[
    list(sfs.k_feature_idx_)]
\end{lstlisting}

\subsection{Embedded: LASSO \zhint{(嵌入法：L1正则化)}}
\zhint{L1正则化会自动把不重要的特征权重压成0。}
\begin{lstlisting}
from sklearn.linear_model import (
    LogisticRegression)
from sklearn.feature_selection import (
    SelectFromModel)
from sklearn.preprocessing import (
    StandardScaler)
scaler = StandardScaler()
scaler.fit(X_train.values)
sel = SelectFromModel(
    LogisticRegression(C=1, penalty='l1'))
sel.fit(scaler.transform(X_train.values),
    y_train)
selected = X_train.columns[sel.get_support()]
X_train = sel.transform(X_train)
X_test = sel.transform(X_test)
\end{lstlisting}

% ============================================================
\section{PCA \zhint{(主成分分析)}}
% ============================================================

\zhint{把很多特征压缩成少数几个"主成分"，保留尽可能多的信息（方差）。}

\keybox{
\zhint{核心思想：找到数据方差最大的方向，投影上去。}\\
\zhint{PCA是无监督的（不看标签），返回最多min(行数,列数)个主成分。}
}

\begin{itemize}[leftmargin=*, nosep, itemsep=0pt]
    \item PC basis = eigenvectors of covariance matrix \zhint{(主成分=协方差矩阵的特征向量)}
    \item Larger eigenvalue = more important PC \zhint{(特征值越大越重要)}
    \item Pick top PCs covering $\geq$ 85\% variance \zhint{(选前几个，覆盖85\%+方差)}
\end{itemize}

\begin{lstlisting}
from sklearn.preprocessing import StandardScaler
from sklearn.decomposition import PCA
from sklearn import datasets
digits = datasets.load_digits()
# Standardize first!
X = StandardScaler().fit_transform(
    digits.data)
# Keep 85% of variance
pca = PCA(n_components=0.85, whiten=True)
X_pca = pca.fit_transform(X)
\end{lstlisting}

% ============================================================
\section{LDA \zhint{(线性判别分析)}}
% ============================================================

\zhint{和PCA类似做降维，但LDA是有监督的——会利用类别标签，让不同类尽量分开。}

\keybox{
\textbf{PCA vs LDA}:\\
\zhint{PCA：无监督，最大化方差（保留信息最多）}\\
\zhint{LDA：有监督，最大化类间距离（分类效果最好）}\\
\zhint{PCA输出最多min(n,p)个PC；LDA输出最多K$-$1个LD（K=类别数）}
}

\begin{lstlisting}
from sklearn.discriminant_analysis import (
    LinearDiscriminantAnalysis as LDA)
from sklearn.preprocessing import StandardScaler
from sklearn.model_selection import (
    train_test_split)
# Split and scale
X_train, X_test, y_train, y_test = \
    train_test_split(X, y,
    test_size=0.2, random_state=0)
sc = StandardScaler()
X_train = sc.fit_transform(X_train)
X_test = sc.transform(X_test)
# LDA: reduce to 1 dimension
lda = LDA(n_components=1)
X_train = lda.fit_transform(X_train, y_train)
X_test = lda.transform(X_test)
\end{lstlisting}
\zhint{注意：LDA的fit需要传y（标签），PCA不需要。}

\subsection{Other Methods \zhint{(其他降维方法)}}
\begin{itemize}[leftmargin=*, nosep, itemsep=0pt]
    \item \textbf{t-SNE} \zhint{(保留局部结构，适合可视化)}
    \item \textbf{UMAP} \zhint{(比t-SNE更快，兼顾全局)}
    \item \textbf{SVD} \zhint{(奇异值分解，PCA的推广)}
\end{itemize}

% ============================================================
\section{Quick Reference \zhint{(速查总结)}}
% ============================================================

\begin{center}
\scriptsize
\begin{tabular}{lp{4cm}}
\toprule
\textbf{Task} \zhint{(任务)} & \textbf{Key Code} \zhint{(关键代码)} \\
\midrule
Split data \zhint{(划分)} & \texttt{train\_test\_split()} \\
Fill NaN \zhint{(填缺失)} & \texttt{SimpleImputer / IterativeImputer} \\
Encode cat. \zhint{(编码)} & \texttt{LabelEncoder / pd.cut / cat.codes} \\
Deduplicate \zhint{(去重)} & \texttt{drop\_duplicates()} \\
Normalize \zhint{(归一化)} & \texttt{MinMaxScaler / StandardScaler} \\
PCA \zhint{(主成分)} & \texttt{PCA(n\_components=0.85)} \\
LDA \zhint{(判别)} & \texttt{LDA(n\_components=1)} \\
Feature sel. \zhint{(特征选择)} & \texttt{SFS / SelectFromModel(LASSO)} \\
\bottomrule
\end{tabular}
\end{center}

\keybox{
\textbf{Preprocessing Pipeline} \zhint{(完整流程)}:\\
Load $\to$ Wrangle $\to$ Impute $\to$ Encode $\to$ Deduplicate $\to$ Split $\to$ Normalize $\to$ Feature Select\\
\zhint{加载 $\to$ 整理 $\to$ 填缺失 $\to$ 编码 $\to$ 去重 $\to$ 划分 $\to$ 归一化 $\to$ 选特征}
}

\end{multicols}
\end{CJK}
\end{document}
